\scanpage{143}

Une courbe elliptique $E$ sur $\mathbb{C}$ est déterminée par
\[
\left\{\begin{aligned}
&\text{a) ``réseau'' }H,\ \deleted{\unsure{\text{le}}}\ \mathbb{Z}\text{-module libre de rang }2\quad\deleted{\unsure{(H=1}}\\
&(H=\pi_1(E,0_E)\simeq H_1(E,\mathbb{Z})=H^1(E,\mathbb{Z})^\vee)\\
&\text{b) Une structure complexe sur }V=H\otimes_{\mathbb{Z}}\mathbb{R}\ \deleted{\unsure{\simeq V|=V}}\ (\simeq\mathfrak{t}_E,\\
&\qquad\text{espace tangent à l'origine de }E).
\end{aligned}\right.
\tag*{(1)}
\]
La donnée a) permet de reconstituer $E$ comme groupe de Lie réel par
\[
E\simeq H\otimes_{\mathbb{Z}}\mathbb{R}\,/\,H
\tag*{(2)}
\]
et la donnée b) \unsure{redonne|redonnent} la structure complexe sur $E$, par passage au quotient.

Quand on est sur une base $S$, espace $\mathbb{C}$-analytique, la donnée a) correspond à un système local. $\mathscr{H}$ sur $S$ -- faisceau en $\mathbb{Z}$-modules loc libre de rang $2$ sur $S$, d'où un fibré vectoriel $\mathbb{R}$-analytique $\mathscr{H}\otimes_{\mathbb{Z}}\mathbb{R}$ sur $S$. La structure b) correspond à une structure complexe \deleted{sur} de ce dernier, en tant que fibré vectoriel \inserted{complexe} de rang $1$, de façon que les sections locales de $\mathscr{H}$ deviennent des sections analytiques complexes. Si \deleted{\unsure{$a$|$e$}} $z_1,z_2$ sont deux sections de $\mathscr{H}$ formant

% archive sheet number: 142
